\section{Discussion}
\label{ch:discussion}

\subsection{Present interpretation of the candidate screen}

The nomination-and-retest screen identified reduced carried-state persistence
as the strongest computational account of the selected cross-task behavioural
pattern against the native baseline. Scaling the carried-state coefficient by
0.95 left ordinary circular accuracy approximately intact while modestly
slowing response settling, increasing vulnerability at longer delays and under
distraction, and selectively impairing 2-back discriminability.

The evidential weight is uneven across tasks. The N-back load contrast was
robust: mean selectivity $0.250\pm0.077$ with a seed-level 95\% interval that
excluded zero and agreement in all ten checkpoints. By contrast, the circular
settling, delay, and distractor contrasts had the predicted mean signs and
majority directional agreement (4/5), but their across-checkpoint intervals all
included zero. In plain terms, the dissertation currently has one solid
behavioural pillar---disproportionate high-load updating impairment---and three
weaker circular hints. The ``complete pattern'' label therefore reflects the
pre-specified majority-direction rule, not equal robustness of every component.

Equation~\ref{eq:persistence} still supplies a coherent model-level reading.
Each update retains slightly less of the preceding hidden state without
simultaneously rescaling incoming drive. Information need not be erased at
once; its stability can weaken cumulatively, so short and simple conditions
remain relatively intact while longer maintenance, competing input, and
repeated updating expose the loss. The conserved time-constant near-misses
support the asymmetry interpretation: when carried-state and drive coefficients
are rescaled together, the joint signature breaks. The present behavioural
analyses do not yet demonstrate a changed attractor or slow manifold; that is
reserved for Section~\ref{ch:results-dynamics}.

\subsection{Relation to the human target profile}

The selective human findings reviewed in Section~\ref{ch:literature} were used
as ordinal targets, not quantitative fits. The robust N-back result aligns in
direction with Barrett-style load dependence~\cite{barrett2018}. The circular
slowing-with-preservation motif is only a directional analogue of
Vollenweider/Yousefi-style RT dominance~\cite{vollenweider1998,yousefi2025}:
settling is externally cued, the mean change is far below one model step, and
the seed-level interval includes zero. Distractor selectivity is likewise only
a weak directional analogue of Carter-style filtering vulnerability
\cite{carter2005}; the circular task tests a related principle, not
multiple-object tracking, and the effect size shrank sharply once distractor
filtering became a trained competence rather than an out-of-distribution
probe. Delay selectivity is also weak after retest, despite cleaner validity.
Thus the human-facing claim must currently be stated as: persistence $0.95$
reproduces the selected load dissociation convincingly and the remaining
circular dissociations only as small majority tendencies.

\subsection{Claim boundary}

``State persistence'' is an operational property of this implemented
recurrence, not a measured neurobiological parameter. The current result shows
computational sufficiency against the native baseline for the descriptive
majority pattern, carried primarily by N-back load selectivity. It does not
show that psilocybin reduces neural persistence, that 5-HT2A activation maps to
$p=0.95$, or that the mechanism is unique. A matched generic-damage comparator
remains the outstanding specificity test
(Section~\ref{ch:results-specificity}); the exploratory Gaussian observation
is only a provisional three-seed N-back hint. Several alternative operators
also reproduced individual components, and cross-operator comparisons remain
uncontrolled for overall clean-task cost.

\subsection{Limitations of the present evidence}

Several limitations constrain inference now. First, post-cue settling is not
literal reaction time because response onset is externally imposed. Second,
circular effect sizes are small and checkpoint-heterogeneous; one checkpoint
reversed the distractor contrast and another slightly reversed delay
selectivity. Third, circular delay selectivity shrank several-fold from the
pilot family to the distractor-trained retest on clean trials alone, so effect
size is family-dependent. Fourth, the descriptive screen used many
operator-strength profiles without a confirmatory inferential family.
Fifth, checkpoint consistency demonstrates robustness across trained models,
not population inference about human participants. Sixth, N-back reuse of three
pilot seeds means that arm is only partly independent of nomination.

\subsection{Next empirical steps}

The immediate next experiment should compare persistence $0.95$ with
matched-cost Gaussian disruption on the present circular and N-back families
(Section~\ref{ch:results-specificity}). If specificity fails on N-back, the
leading candidate should be demoted regardless of the descriptive screen. If
specificity holds on N-back but circular contrasts remain weak, the
dissertation story should be rebalanced around load selectivity as the primary
result, with circular outcomes as secondary directional support. Only after
that comparison should persistence be densified around $0.95$ and delay-period
representational drift quantified (Section~\ref{ch:results-dynamics}).
